\section{Security Hardening}
\label{sec:analysis-security}

\subsection{RQ2.1: Service Identity and Mutual TLS}
\label{sec:analysis-rq21}

The results in \cref{sec:rq21} support four bounded findings.

\paragraph{Finding 1: the mTLS/AuthZ hardening bundle adds measurable but bounded overhead to the selected chain.}
For the primary \texttt{chain\_2svc} topology, the service-mesh hardening
layer adds 3.488~ms, or 4.62\%, to mean end-to-end latency
(\cref{tab:rq21_latency}). This is a measurable cost, but it is bounded
relative to the approximately 75.6~ms plain AKS-with-mesh-tainted baseline. The result is consistent with service-mesh literature: sidecars
add critical-path work, but the visible percentage overhead depends on the
workload and on how much application compute the proxy cost is amortised
over~\cite{meshinsight_socc23}.

\paragraph{Finding 2: security overhead grows with protected chain depth.}
The \texttt{chain\_5svc} stress run increases the mean hardening overhead to
8.837~ms (10.74\%). This is approximately two and a half times the absolute
\texttt{chain\_2svc} overhead. The larger cost is expected: the stress
topology contains more injected proxies, more authorization policies, and
more protected inter-service boundaries. This finding reinforces the RQ1
result that service-count decisions and security decisions cannot be
separated. A chain that is tolerable without communication hardening may
become less attractive once every boundary also carries identity, encryption,
proxying, and policy-enforcement cost.

\paragraph{Finding 3: the result is lower than some published Istio overheads, but not in conflict with them.}
The observed 4.62\% primary overhead is much smaller than the high-load Istio
penalties reported in broader service-mesh
benchmarks~\cite{bremlerbarr2025mtls}. This does not contradict that
literature. Those studies use different services, protocols, load levels,
tail-latency metrics, mesh configurations, and deployment environments. RQ2.1
uses a CPU-bound ResNet-18 inference chain with substantial per-request model
computation, same-node placement, and a specific managed-Istio AKS
configuration. The correct comparison is therefore qualitative: both this
thesis and the literature show that sidecar-mediated mTLS has non-zero latency
and resource cost, while the absolute magnitude is workload-specific.

\paragraph{Finding 4: RQ2.1 hardens inter-service communication, not the entire trust model.}
The security validation shows that the intended service-account-based
communication policy was enforced for the protected downstream segments:
valid full-chain requests succeeded, while direct non-mesh probes to protected
downstream services were blocked. This is meaningful communication-level
hardening. However, it does not remove trust in the service-mesh control
plane, does not protect against all Kubernetes administrator threats, and
does not provide confidential execution for model code or activations. Those
limits matter because prior work identifies the mesh control plane and local
CA as security-critical trust points~\cite{poudel2025mazu}. RQ2.1 should
therefore be claimed as a measured first hardening layer, not as a complete
security solution.

\paragraph{Answer to RQ2.1.}
Under the paired AKS RQ2.1 benchmark, adding service identity, managed-Istio
mTLS, and inter-service authorization policy introduces measurable but
bounded overhead for the selected \texttt{chain\_2svc} deployment. Mean
latency increases from 75.557~ms in the plain condition to 79.045~ms with
the mTLS/AuthZ hardening bundle, an overhead of 3.488~ms or 4.62\%; p95
latency increases by 3.623~ms. Security validation passed in every mTLS
pass, including full-chain positive probes and direct-denial probes against
protected downstream segments. The additional \texttt{chain\_5svc} stress run
shows that the same hardening approach becomes more expensive as service
depth increases: mean overhead rises to 8.837~ms or 10.74\%. RQ2.1 therefore
supports a cautious conclusion: service identity and managed-Istio mTLS/AuthZ
are feasible for the selected Azure chain at a modest latency cost, but the
cost is not negligible, grows with the number of protected service
boundaries, and comes with operational complexity from sidecars, policies,
resource requests, and readiness delay. The follow-up ablation indicates that
this cost is dominated by entering the managed-Istio sidecar data path rather
than by strict mTLS enforcement or AuthorizationPolicy evaluation alone.

\subsection{RQ2.2: Confidential VM Execution}
\label{sec:analysis-rq22}

The results in \cref{sec:rq22} support three bounded findings.

\paragraph{Finding 1: VM-level confidential execution adds measurable overhead on top of mTLS.}
The RQ2.2 paired run shows positive confidential-execution overhead relative
to a newly collected standard baseline. This matters because the baseline
already includes the RQ2.1 mTLS/AuthZ hardening. The measured delta therefore
represents the additional cost of placing the downstream service on
confidential VM infrastructure, not the cost of introducing service-mesh
communication security.

\paragraph{Finding 2: selective TEE has a modest measured cost on the protected segment.}
Placing only \texttt{service2} on \texttt{Standard\_DC8as\_v5} (AMD SEV-SNP)
increases mean latency by 1.854~ms (+3.17\%), median latency by 1.467~ms
(+2.52\%), and p95 latency by 3.601~ms (+5.98\%) relative to the matched
standard-node baseline (\cref{tab:rq22_service2_results}). This is the most
direct implementation of the original RQ2.2 design: \texttt{service2} is the
downstream protected component, receives intermediate activations from
\texttt{service1}, and is the service guarded by the strict mTLS and
authorisation policy in the RQ2.1 contract. He~et~al.\
\cite{he2021_attacking_collaborative} motivate protecting that intermediate
activation beyond the in-transit mTLS layer, although they target a
malicious-remote-service threat model that differs from the
infrastructure-level adversary SEV-SNP defends against; selectively placing
\texttt{service2} on confidential infrastructure therefore gives the thesis a
narrow, defensible hardening point with a modest measured cost.

\paragraph{Finding 3: the security claim is stronger than RQ2.1 but still bounded.}
RQ2.1 protected the service-to-service communication path. RQ2.2 adds a
VM-level runtime protection boundary against host-level access for the
confidentially placed service. This is a stronger security posture, but it is
not complete application-level isolation. The guest OS, application process,
model runtime, and local sidecar remain trusted parts of the confidential VM
environment, and plaintext exists inside the endpoint boundary after mTLS
termination. Recent work additionally shows that confidential VMs retain
residual attack surfaces beyond memory encryption, including malicious
interrupt injection from the untrusted
hypervisor~\cite{schluter2024_heckler}; this thesis does not experimentally
evaluate such attacks. The correct claim is therefore that RQ2.2 adds
VM-level confidential execution to the already communication-secured Azure
inference chain, not that it removes every trusted component or eliminates
all possible data exposure.

\paragraph{Answer to RQ2.2.}
RQ2.2 shows that AMD SEV-SNP VM-level confidential execution is feasible for
the communication-secured AKS inference chain and introduces a measurable but
bounded additional overhead. With only the protected downstream service moved
to a confidential node (\texttt{service2\_only} scope), mean latency
increases by 1.854~ms, or 3.17\%, and p95 latency increases by 3.601~ms,
relative to a newly paired standard-node mTLS baseline. Extending the
confidential-execution boundary to both services in the chain (full
two-service TEE) is out of the thesis-facing scope of RQ2.2 per the evidence
manifest and is treated as future work. The interpretation remains bounded by
the SEV-SNP threat model and by the sidecar-based service-mesh architecture:
RQ2.2 evaluates VM-level confidential execution on Azure for the downstream
segment, not application-level enclave isolation or complete end-to-end
secrecy.

\subsection{RQ2.3: Trade-off Synthesis}
\label{sec:analysis-rq23}

\Cref{sec:rq23} shows that the evaluated security mechanisms form a
cumulative hardening gradient rather than a single binary secure/insecure
choice. The RQ1.5 plain AKS deployment establishes the performance baseline.
RQ2.1 adds communication-level security through mTLS, service identity, and
authorization policy. RQ2.2 then adds VM-level confidential execution for the
downstream protected service (\texttt{service2\_only}). Each step increases
the protection scope, but each step also introduces additional latency,
resource, or operational cost. Extending the confidential-execution boundary
to both services (full two-service TEE) is treated as future work and is not
evaluated as a measured trade-off point in this thesis.

The main engineering implication is that security hardening should be scoped
to the threat model. Applying every available mechanism everywhere is not
always the best default. In the evaluated deployment, mTLS/AuthZ is the most
suitable general-purpose security baseline because it protects the
inter-service path at a modest cost for \texttt{chain\_2svc}. Selective TEE
on the downstream service is the next available hardening step when that
segment is the primary sensitivity
boundary~\cite{he2021_attacking_collaborative}, while its security claim
remains bounded by the SEV-SNP threat
model~\cite{schluter2024_heckler}.

The results also show that performance and security cannot be evaluated
independently. RQ1 established that service boundaries add measurable
overhead and that this overhead depends on boundary placement and service
count. RQ2 shows that security mechanisms compound this cost. The mTLS/AuthZ
hardening overhead grows when more service boundaries are protected, and
selective confidential execution adds further overhead on top. The final
deployment decision therefore depends on three linked factors: the chosen
model partition, the number of service boundaries, and the required
protection scope.

\paragraph{Answer to RQ2.3.}
RQ2.3 shows that the best hardening level depends on the assumed threat
model and the acceptable performance cost. Plain AKS is the fastest
evaluated configuration, but it provides the weakest protection in the
security comparison. mTLS with service identity and authorization policy is
the best general-purpose balance when the main goal is authenticated and
encrypted service-to-service communication; in the selected
\texttt{chain\_2svc} deployment it adds 3.488~ms, or 4.62\%, mean latency
overhead. Selective TEE on \texttt{service2\_only} is the preferred option
evaluated in this thesis when the downstream inference segment is the main
infrastructure-level sensitivity boundary, adding 1.854~ms, or 3.17\%, over
its paired standard-node mTLS baseline. Extending the confidential-execution
boundary to both services (full two-service TEE) is treated as future work
and is not evaluated here as a measured trade-off point.

Overall, the evaluated configurations support a cumulative
security-performance trade-off: stronger protection is feasible for the
selected AKS inference chain, but each additional protection layer adds
measurable cost. The most defensible default is mTLS/AuthZ for
communication-level protection, with selective confidential execution on the
downstream service added only when the threat model requires runtime
protection against host-level exposure, and bounded by the SEV-SNP
attack-surface limitations discussed
above~\cite{schluter2024_heckler}.

\subsection{Synthesis: Security Hardening}
\label{sec:analysis-security-synthesis}

The security stages show that hardening is feasible for the selected AKS
deployment, but it is not free. RQ2.1 adds managed-Istio mTLS, service
identity, and authorization policy. For the selected \texttt{chain\_2svc}
topology, the measured mean overhead is 3.488~ms, or 4.62\%; for the
\texttt{chain\_5svc} stress topology, the overhead grows to 8.837~ms, or
10.74\% (\cref{tab:rq21_latency,tab:rq21_paired_deltas}). This is consistent
with service-mesh studies showing that sidecar-based mTLS adds latency and
resource cost through additional critical-path
processing~\cite{meshinsight_socc23,bremlerbarr2025mtls}.

RQ2.2 adds VM-level confidential execution on an AMD SEV-SNP node for the
downstream protected service while keeping the mTLS contract fixed.
Selective confidential execution for \texttt{service2\_only} adds 1.854~ms,
or 3.17\%, over its paired standard-node mTLS baseline. The result is
consistent with the confidential-computing literature's expectation that
CVM overhead is workload- and
scope-dependent~\cite{yan2023_cvm_overheads,misono2024_cvm_explained,
mo2024_ml_confidential_computing}, and the security claim is bounded by
SEV-SNP's residual attack
surfaces~\cite{schluter2024_heckler}.
